% Source presentations for register row N-c-015.

\subsection{Finite trees in \emph{Tree Automata Techniques and Applications}}
\label{sec:p09-presentation-n-c-015}

The source is the online monograph release dated 18 November 2008; the register's
2007 date refers to the evolving work rather than a distinct fixed book
edition.  The release explicitly treats finite trees and excludes infinite-tree
automata \cite[pp.~11--13]{ComonEtAl2007}.  Its automata therefore receive
finite terms, not Rabin's entire infinite labeled tree.

A finite ranked alphabet \(F\) assigns an arity to each of its symbols.  Write
\(T(F)\) for its set of ground terms, formed recursively: every nullary symbol
is a term, and
\(f(t_1,\ldots,t_n)\) is a term exactly when \(f\) has arity \(n\) and each
\(t_i\) is already a term.  Their node diagrams are the finite ordered
\(F\)-trees.  TATA writes a nondeterministic bottom-up automaton as
\[
 \mathcal A=(Q,F,Q_f,\Delta),
\]
where \(Q\) is the finite state set, \(Q_f\subseteq Q\) the final states, and
\(\Delta\) contains rules
\(f(q_1,\ldots,q_n)\to q\) for \(n\)-ary \(f\); several rules may have
the same symbol and child-state tuple.  A run labels each tree node
bottom-up so that its label and its children's states match one transition; it
accepts when the root state lies in \(Q_f\)
\cite[\S1.1]{ComonEtAl2007}.  A language is \emph{recognizable} when some
such automaton accepts exactly its members.  This tuple is the source notation
corresponding to Definition~\ref{def:p09-ranked-tree-automaton}: TATA's
alphabet \(F\) is the handbook's \(\Sigma\), so \(T(F)\) is the handbook's
\(T_\Sigma\); TATA's final-state set \(Q_f\) is the handbook's \(F\), and TATA
lists \(Q_f\) before \(\Delta\).

For a small run, take nullary symbols \(a,b\), a binary symbol \(f\), rules
\(a\to q_a\), \(b\to q_b\), and
\(f(q_a,q_b)\to q_f\), with \(q_f\) final.  The unique displayed bottom-up
assignment accepts \(f(a,b)\).  If the last rule is absent, the same ground
term has no run.  This example makes clear that acceptance concerns the root of
a finite term rather than recurrence along an infinite path.

\paragraph{Subset determinization.}
\label{sec:p09-finite-tree-determinization}
For a term \(t\), let \(\operatorname{Reach}_{\mathcal A}(t)\subseteq Q\)
be the possible root states of all runs on \(t\), without imposing
the final-state condition.
The deterministic automaton has state set \(\mathcal P(Q)\), final
states
\[
 \mathcal F=\{S\subseteq Q:S\cap Q_f\ne\varnothing\},
\]
and, for each symbol \(f\) of arity \(m\), the total transition function
\[
 \delta_f(S_1,\ldots,S_m)
 =\{q\in Q:\exists q_1\in S_1\cdots\exists q_m\in S_m\
                    (f(q_1,\ldots,q_m)\to q)\in\Delta\}.
\]
For \(m=0\), this is the set of targets of all nullary \(f\)-rules.
The empty subset is retained, so the transition is defined even
when there is no compatible original rule.

Induction on \(t=f(t_1,\ldots,t_m)\) proves that the unique
deterministic root state is \(\operatorname{Reach}_{\mathcal A}(t)\).
A run on \(t\) restricts to runs on its child terms and one matching
root rule, giving membership in the displayed set.
Conversely, the chosen child runs and that root rule combine into
a run on \(t\).  For a nullary term the assertion is precisely
the definition of \(\delta_f\).
The final-state condition therefore gives exactly \(L(\mathcal A)\).
There are at most \(2^{|Q|}\) states.

One may retain only the subsets realized by ground terms.
Starting with nullary transition values and repeatedly applying all
symbol operations to tuples of already obtained subsets computes
them.  Every generated subset has a term witness, and induction on
terms gives the converse.  If the ranked alphabet has no nullary
symbol, \(T(F)\) is empty; the full subset construction still defines
a complete automaton, and there is no input term to accept
\cite[Theorem~1.1.9]{ComonEtAl2007}.

\paragraph{Boolean operations on tree languages.}
\label{sec:p09-finite-tree-boolean}
Complement the deterministic complete automaton by replacing
\(\mathcal F\) with \(\mathcal P(Q)\setminus\mathcal F\).
Every ground term has exactly one run, so its root belongs to exactly
one of these final-state sets.  The result is the complement relative
to \(T(F)\).
The totality and uniqueness of the run are used here: flipping final
states in a nondeterministic automaton can retain an accepted term
through a different run, and an absent run in a partial automaton
cannot become accepting merely by changing its final states.

For union, first make the state sets of two automata over the same
ranked alphabet disjoint.  Take their disjoint union of states, rules,
and final states.  Every transition keeps all its child states and
its target in one component.  A run on the connected input tree
therefore lies entirely in one component, and acceptance is precisely
acceptance by either original automaton.

For intersection, use states \(Q_1\times Q_2\), final states
\(Q_{1,f}\times Q_{2,f}\), and rules
\[
 f((q_1,r_1),\ldots,(q_m,r_m))\to(q,r)
\]
exactly when both \(f(q_1,\ldots,q_m)\to q\) and
\(f(r_1,\ldots,r_m)\to r\) are original rules.
Projection of a run gives a run of each component, and pairing
two runs on the same tree gives a product run.
The root is final exactly when both original roots are final.
These are the effective Boolean closures of
\cite[Theorem~1.3.1]{ComonEtAl2007}.

\paragraph{Membership, emptiness, and finite witnesses.}
\label{sec:p09-finite-tree-emptiness}
Membership of a supplied finite term is decided by computing its
reachable-state subsets upward and checking the root.
For emptiness, start with the states reached by nullary rules and
repeatedly mark the target of any rule whose child states are all marked.
Induction on the marking stages shows that every marked state is
the root of a run on some ground term.
Conversely, induction on term height shows that every possible root
state is eventually marked.
Thus the language is nonempty exactly when some final state is marked.
At most \(|Q|\) new-state markings occur.

For the linear-time implementation, give each rule a counter equal
to its arity and index each child-state occurrence by its rule.
Enable nullary rules immediately.
Whenever a state is first marked, process each of its indexed
occurrences and decrement the corresponding rule counter.
When a counter reaches zero, mark the rule's target if it is new.
Occurrences are counted with multiplicity: a rule
\(f(q,q)\to r\) has two indexed occurrences of \(q\), and both are
processed when \(q\) is marked.
Each state enters the queue once, and each child occurrence is
processed once.  With indexed state and rule identifiers, the running
time is
\[
 O\!\left(|Q|+|\Delta|+\sum_{\rho\in\Delta}
                                 \operatorname{arity}(\rho)\right),
\]
linear in the explicit transition representation
\cite[Theorem~1.7.4]{ComonEtAl2007}.

Recording the rule that first marks each state also gives a finite
accepting witness when one exists.
Its child states were marked earlier, so these recorded dependencies
form an acyclic finite graph.  Unfolding the graph gives a ground
term and an accepting run.
The graph representation can share repeated subterms; writing out
the full tree may require more space than deciding nonemptiness.

Combining complement and intersection with emptiness decides inclusion:
\[
 L(\mathcal A)\subseteq L(\mathcal B)
 \quad\Longleftrightarrow\quad
 L(\mathcal A)\cap\bigl(T(F)\setminus L(\mathcal B)\bigr)=\varnothing.
\]
Two inclusion decisions give equivalence; emptiness of the complement
gives universality.  These constructions use only finite state sets
and finite transition tables.

For \(k\geq1\), let \(\epsilon\) denote the empty word in
\(\{1,\ldots,k\}^{*}\).  WSkS is weak monadic second-order logic over this
set of words: first-order variables denote words, the successor
operation appends one letter, and every set variable denotes a finite set of
words.  Proposition~3.3.1 replaces each first-order variable by a singleton set
variable.  The resulting restricted syntax has atomic formulas
\[
 X\subseteq Y,\qquad \operatorname{Sing}(X),\qquad X=Yi,\qquad X=\epsilon,
\]
where \(\operatorname{Sing}(X)\) says that \(X\) has exactly one element and
\(X=Yi\) says that the unique word in \(X\) is the \(i\)-successor of the
unique word in \(Y\).  The source atom \(X=\epsilon\) is the eliminable
abbreviation \(X=\{\epsilon\}\).  This removes first-order variables without
changing the defined assignments.

The translation uses a fresh set variable for each individual variable.
An existential individual quantifier becomes
\(\exists X\,(\operatorname{Sing}(X)\land\cdots)\), and a universal
one becomes \(\forall X\,(\operatorname{Sing}(X)\to\cdots)\).
After the recursive translation, conjoin singleton guards for the
free individual-variable tracks.
The bijection \(w\mapsto\{w\}\) between words and singleton sets
proves that these clauses preserve their assignment domains.

For \(n\geq1\), a tuple \((S_1,\ldots,S_n)\) of finite sets is encoded by one finite tree over
\(\{0,1,\bot\}^{n}\).  At a position \(p\), coordinate \(i\) is \(1\) when
\(p\in S_i\), \(0\) when \(p\notin S_i\) but extends to a word in \(S_i\),
and \(\bot\) otherwise.  A symbol containing \(0\) or \(1\) has arity \(k\),
whereas \(\bot^{n}\) is nullary.
For precision, let \(K=\{1,\ldots,k\}\), put \(U=\bigcup_i S_i\),
and let \(\operatorname{Pref}(U)\) contain all prefixes of members of \(U\).
The code has exactly the position domain
\[
 D(\bar S)=\{\epsilon\}\cup
                \{pi:p\in\operatorname{Pref}(U),\ i\in K\}.
\]
Thus each prefix supporting a set member has all \(k\) children;
children with no support are nullary padding leaves.
The prefix set and this domain are finite because \(U\) is finite.
When all \(S_i\) are empty, the code is the one-node tree \(\bot^n\).
Write \(\operatorname{code}(\bar S)\) for the resulting term.
Decoding selects precisely the positions with a \(1\) in each track.

The source's Example~3.3.3, with
\(k=2\), \(S_1=\{\epsilon,11\}\), \(S_2=\varnothing\), and
\(S_3=\{11,22\}\), labels the root by \((1,\bot,0)\), positions \(1\) and
\(2\) by \((0,\bot,0)\) and \((\bot,\bot,0)\), positions \(11\) and \(22\)
by \((1,\bot,1)\) and \((\bot,\bot,1)\).
Its six padding positions are exactly
\(12,21,111,112,221,222\), each labelled \((\bot,\bot,\bot)\).
There are eleven nodes in total.

\paragraph{Recognizing exactly the canonical codes.}
\label{sec:p09-wsks-canonical}
For \(n\geq1\), write \(\Sigma_n=\{0,1,\bot\}^n\) with the ranks
just defined, and \(\bot_n=\bot^n\).
Let \(W_n\subseteq T(\Sigma_n)\) be the set of canonical codes.
A partial deterministic bottom-up checker has states
\(H\subseteq\{1,\ldots,n\}\), recording which tracks have a
membership mark somewhere in the current subtree.
Its nullary rule is \(\bot_n\to\varnothing\).
For an internal label \(\alpha\) and child states \(H_1,\ldots,H_k\),
put \(H=\bigcup_{\ell=1}^k H_\ell\) and permit a transition precisely when
\[
\begin{array}{ll}
 \alpha_j=\bot&\Longrightarrow j\notin H,\\
 \alpha_j=0&\Longrightarrow j\in H,\\
 \alpha_j=1&\text{has no additional child condition}
\end{array}
\qquad(1\leq j\leq n).
\]
The target is \(H\cup\{j:\alpha_j=1\}\), and every checker state
is final.  An invalid local combination has no transition.
Induction on subtrees proves the support interpretation of the states.
The three conditions say exactly that bottom forbids all membership
at or below the node, zero requires a strict descendant member,
and one marks membership at the node itself.
Every internal node consequently supports a member in some track.
Together with the ranks, this forces the displayed code domain
and all its labels.  Conversely every canonical code satisfies these
conditions, so the checker recognizes exactly \(W_n\).

For example, \(00(\bot\bot,\bot\bot)\) is not a canonical binary
two-track code: both zeros assert support that neither child has.
The checker rejects it.
For the empty variable list, set
\(\Sigma_0=\{\top\}\), with \(\top\) nullary, and \(W_0=\{\top\}\);
this single tree encodes the empty tuple.
Throughout the formula construction we maintain
\[
 L(\mathcal A_\varphi)
   =\{\operatorname{code}(\bar S):\bar S\models\varphi\}
   \subseteq W_n,
\]
where the coordinates are the distinct free set variables in a fixed order.

This coding also supplies the exact bridge to finite-tree monadic logic.  For
a finite ranked alphabet \(F\), choose \(k\geq1\) at least as large as every
arity occurring in \(F\).  A finite ranked tree is represented by its finite
position set \(D\) and the
subsets \(X_f\subseteq D\) of positions carrying each symbol \(f\); the
source's arity clauses say that these data form one ground term.  In that
coding the label predicate for \(f\) becomes \(x\in X_f\), and the
\(i\)-th-child relation becomes \(y=xi\).  A monadic formula over the finite
tree becomes a WSkS formula by restricting every element quantifier to \(D\)
and every set quantifier to subsets of \(D\).
Conversely, the state sets in Lemma~3.3.6 below partition precisely \(D\).
Thus the ambient prefix tree is used to code a finite input, not substituted
for that input \cite[pp.~89--95]{ComonEtAl2007}.

Lemma~3.3.4 constructs an automaton from a restricted WSkS formula.  In the
binary instance \(k=2\), its atomic construction for \(X\subseteq Y\), for
example, has one final state \(q\) and the rules
\[
 \bot\bot\to q,\quad 00(q,q)\to q,\quad 01(q,q)\to q,\quad
 11(q,q)\to q,\quad \bot0(q,q)\to q,\quad \bot1(q,q)\to q.
\]
Intersect this automaton with \(W_2\).
On a canonical code of an inclusion, a nonempty \(X\)-support
cannot face an empty \(Y\)-support, and every \(X\)-membership mark
also marks \(Y\).  Exactly the displayed labels are permitted.
Conversely, the absence of \(10\) and \(1\bot\) at every node
ensures that every member of \(X\) belongs to \(Y\).
For general \(k\), use the same five internal labels with \(k\)
children in state \(q\).

The remaining atomic constructions are equally finite.
For \(\operatorname{Sing}(X)\), use states \(q,q'\), with \(q'\)
the only final state, and rules
\[
 \bot\to q,\qquad 1(q,\ldots,q)\to q',
\]
together with, for each child position \(i\), a \(0\)-rule with
child \(i\) in state \(q'\), all other children in \(q\), and target \(q'\).
A \(q'\)-run therefore follows one zero-labelled prefix path to
one membership node; all other subtrees are empty.
These are precisely singleton codes.
For \(X=\epsilon\), retain just the two displayed rules, so the
membership node must be the root.

For distinct variables in the order \((X,Y)\), the atom \(X=Yi\)
uses states \(q,q',q''\), with only \(q''\) final, and rules
\[
 \bot\bot\to q,\qquad 1\bot(q,\ldots,q)\to q'.
\]
Add one \(01\)-rule whose \(i\)-th child is in \(q'\), whose other
children are in \(q\), and whose target is \(q''\).
Finally add the \(00\)-rules with exactly one child in \(q''\),
all others in \(q\), and target \(q''\).
The \(01\)-node is the unique \(Y\)-member; its \(i\)-th child is
the unique \(X\)-member.  The \(00\)-rules allow their common prefix
above that pair.  This proves the required singleton-successor
condition.  Repeated variables are handled separately:
\(X\subseteq X\) accepts \(W_1\), whereas \(X=Xi\) accepts nothing
because a finite word is not its own proper successor.
These constructions are the general-\(k\) forms of the source's
binary transition tables \cite[Figures~3.8--3.11]{ComonEtAl2007}.
Permuting coordinates puts each atomic automaton into the chosen
variable order; the operation preserves every symbol's arity.
The inverse-image construction below supplies any additional free tracks.

\paragraph{Projection with subtree pruning.}
\label{sec:p09-wsks-pruning}
For \(n>1\), let \(d_j\) delete coordinate \(j\) from a label,
and put \(m=n-1\).
Define a tree map \(h_j\) by
\[
 h_j(\bot_n)=\bot_m,
\]
and
\[
 h_j(\alpha(t_1,\ldots,t_k))=
 \begin{cases}
  \bot_m,&d_j(\alpha)=\bot_m,\\
  d_j(\alpha)(h_j(t_1),\ldots,h_j(t_k)),
       &d_j(\alpha)\ne\bot_m.
 \end{cases}
\]
The first case deletes the entire child family because the projected
symbol is nullary.
For canonical input, it occurs exactly where no surviving track
has a member at or below the node.  Otherwise a surviving membership
keeps the node internal, and the same calculation applies recursively
to its children.  Hence
\[
 h_j(\operatorname{code}(S_1,\ldots,S_n))
   =\operatorname{code}(S_1,\ldots,S_{j-1},S_{j+1},\ldots,S_n).
\]
A surviving zero still has a descendant one in its own track, so
its label remains zero.  This is the source's linear tree
homomorphism specialized to these ranked codes
\cite[Proposition~3.2.12]{ComonEtAl2007}.

Here is the automaton image construction, including the collapsed case.
Let \(\mathcal A\) recognize a language contained in \(W_n\).
First take the product of \(\mathcal A\) with the partial checker
for \(W_n\).  Its accepted language is unchanged, and every tree
admitting a run is now canonical.
Let \(E\) be the set of states reachable by trees whose surviving
coordinates are bottom at every node.
Compute \(E\) by the emptiness marking construction restricted to
all nullary \(\bot_n\)-rules and the internal labels \(\alpha\) satisfying
\(d_j(\alpha)=\bot_m\).
The resulting states have finite witness subtrees supported entirely
by the deleted track.

Keep the state set and final states, and use the following projected rules:
\[
\begin{array}{ll}
 \bot_m\to q&\text{for every }q\in E,\\
 d_j(\alpha)(q_1,\ldots,q_k)\to q
   &\text{for every old }\alpha(q_1,\ldots,q_k)\to q\\
   &\text{with }d_j(\alpha)\ne\bot_m.
\end{array}
\]
An original accepting run projects by retaining its state at every
remaining node.  At a collapsed subtree the root state is in \(E\),
so the new nullary rule is available.
Conversely, replace each new nullary leaf of a projected accepting
run by a finite subtree witnessing its state in \(E\).
At each retained internal node use an old transition witnessing
the new rule.  Induction up the finite tree produces an original
accepting run whose image is the projected tree.
That original tree is canonical, so the displayed coding identity
proves that the image represents exactly the surviving assignment.

For example, deleting \(Y\) from a code with \(X=\varnothing\)
and \(Y=\{1\}\) collapses its root to the one-node empty-\(X\) code.
The old root's surviving label is bottom, so keeping its children
would violate the required rank.

\paragraph{Aligning free variables and interpreting connectives.}
\label{sec:p09-wsks-logical-induction}
Suppose an automaton uses only some of the variables in a larger list,
with list lengths \(m>n\geq1\).
Let \(h\) drop all other tracks with the pruning rule just proved.
Its language on the larger list is
\[
 W_m\cap h^{-1}(L(\mathcal A)).
\]
This inverse image has an explicit finite construction.
Determinize and complete the smaller automaton, with transition
functions \(\delta\).  At a larger-alphabet node, let \(\beta\)
be the vector of retained coordinates.
If \(\beta=\bot_n\), assign the fixed state \(\delta_{\bot_n}\),
ignoring the child states.  Otherwise apply \(\delta_\beta\) to
the child states.  A larger padding leaf also receives
\(\delta_{\bot_n}\).
The resulting state is the smaller automaton's state on \(h(t)\),
by induction on \(t\).
Totality ensures that even a child whose image is rejected has a
state and cannot obstruct a later collapse above it.
Intersecting with \(W_m\) completes the construction.

For a disjunction, lift both operand automata to the union of their
distinct free-variable lists and take union.
For negation use relative complement
\[
 L(\mathcal A_{\neg\varphi})=W_n\setminus L(\mathcal A_\varphi).
\]
For conjunction, apply the same free-variable alignment and then
take intersection.
For \(\exists X\,\varphi\), take the pruned image deleting \(X\)'s
track.  Every preimage is a finite-set witness, and every finite-set
witness has its canonical preimage, so the operation has exactly
the existential semantics.  If \(X\) is not free in \(\varphi\),
the quantifier is vacuous and the automaton is unchanged.
Bound variables are renamed first so their track identities are unambiguous.
Together with the atomic constructions, these cases prove the
language invariant by induction on restricted WSkS formulas.

At zero free variables, \(\Sigma_0\) has just the nullary symbol
\(\top\), and \(W_0=\{\top\}\).
Deleting the last track sends every code to this one nullary tree.
The set \(E\) in the image construction is then the set of all
reachable states, and the rules are \(\top\to q\) for \(q\in E\).
Thus \(\top\) is accepted exactly when some old accepting code exists,
including witnesses of arbitrarily large finite depth.
After determinizing and completing a zero-track automaton, lifting
uses its fixed state \(\delta_\top\) at every node and then intersects
with the larger canonical-code language.
Complement at level zero exchanges \(\varnothing\) and \(\{\top\}\).

Conversely, Lemma~3.3.6 assigns a set variable \(Y_q\) to the positions bearing
each automaton state \(q\), requires those sets to partition the encoded tree,
and requires the root to satisfy
\(\bigvee_{q\in Q_f}\epsilon\in Y_q\).  If \(X_f\) is the set of positions
bearing symbol \(f\), its transition condition is the source's local clause
\[
 \forall x\Bigl((x\in X_f\land x\in Y_q)\Rightarrow
   \bigvee_{f(q_1,\ldots,q_s)\to q\in\Delta}
      \bigwedge_{i=1}^{s}xi\in Y_{q_i}\Bigr),
\]
conjoined for every symbol \(f\) and state \(q\).  These clauses say exactly
that the quantified state partition is an accepting bottom-up run
\cite[\S3.3]{ComonEtAl2007}.

Lemmas~3.3.4 and~3.3.6 supply the two coded directions: a formula
gives an automaton on assignment codes, and an accepting run is
described by an existential state partition.
Theorem~3.3.8 then obtains WSkS decidability.  For a
closed sentence the track alphabet is the singleton nullary alphabet
\(\{\top\}\), so truth in the fixed WSkS structure reduces to
whether a rule \(\top\to q_f\) exists for some final state \(q_f\).  This is
the detailed source route for Theorem~\ref{thm:p09-finite-tree-automata}.

The restriction is real.  For the ranked alphabet
\(F=\{a^{(0)},f^{(2)}\}\), Exercise~1.4 records that
\(L=\{f(t,t):t\in T(F)\}\) is not recognizable.  By Theorem~1.1.9, a
recognizing automaton could be taken deterministic.  Every term \(t\) would
have to reach a state, since otherwise the required term \(f(t,t)\) could not
be accepted.  Distinct terms \(t,u\) could not reach the same state: otherwise
the identical transition on the child-state pair would give \(f(t,t)\) and
\(f(t,u)\) the same root state, although the first lies in \(L\) and the second
does not.  Thus the infinitely many terms in \(T(F)\) would require infinitely
many states, a contradiction.

If set variables range over arbitrary, not necessarily finite, subsets of the
infinite successor tree, the source calls the resulting logic full SkS.  Its
decision procedure requires automata that assign states from the root downward
on infinite inputs and a different complementation argument; TATA states that
boundary immediately
after Theorem~3.3.8.  On an actually finite input structure, weak and full set
quantification coincide because every subset is finite.  That elementary
observation does not turn a finite-tree automaton into Rabin's automaton on an
entire infinite labeled tree.

\paragraph{Thatcher--Wright source boundary.}
The registered 1968 article is the historical source named by TATA.  Its
authoritative abstract says only that finite algebras replace conventional
finite automata and announces closure results, a generalized Kleene theory,
and an application to a positive second-order decision problem.  Because the
technical body is not part of the available source basis, that abstract-level
phrase is not used as a handbook definition.  The paper's algebraic signature,
recognition relation, exact logical language, theorem statements, examples, and
proof are not reconstructed by assigning TATA's modern tuple and WSkS proof to the older
paper \cite{ThatcherWright1968}.  Historical provenance and the detailed
modern presentation therefore remain separate.  A same-title 1966 IBM
Research Paper RC-1713 is a bibliographic lead only, not technical evidence.

\paragraph{Finite-domain MSO and the successor-tree coding.}
The finite-tree equivalence in Theorem~\ref{thm:p09-finite-tree-automata}
has a complete proof in Annex~\ref{annex:p9.trees.1}, using assignment
bits on the nodes of a fixed finite input tree.
Erasing those bits preserves the input tree's ranks and shape.
The WSkS coding here instead represents finite sets in the infinite
successor tree, and its projection includes the proved subtree pruning.
Both constructions are explicit at the input type where they are used.
